\chapter{Geodesics, Uniform Acceleration, and the Newtonian Limit}

The immediate problem is gravity. More precisely, we want the spacetime metric that describes a nearly uniform gravitational field, such as the field in a small laboratory near the surface of the Earth. The shortest route to that metric turns out not to be a direct one. We must first make three geometric ideas operational: covariant differentiation, parallel transport, and geodesic motion.

These ideas depend on one another. Covariant differentiation tells us how to compare vectors at neighboring points. Parallel transport tells us how to carry a vector without turning it. A geodesic is obtained when the vector being carried is the tangent to the path itself. Once the path lies in spacetime rather than in space, that same geodesic equation becomes the equation of motion of a freely falling particle.

The second half of the lecture then makes a striking move. We remain in flat Minkowski spacetime, but describe it with coordinates adapted to a rigidly accelerated family of observers. In those coordinates the metric looks gravitational. Near one observer it reproduces the potential of a uniform field, and its geodesics obey the Newtonian equation
\begin{equation}
\frac{d^2y}{d\tau^2}=-g.
\end{equation}
This is the equivalence principle written as a calculation.

\section{The Covariant Derivative Recalled}

At a chosen point \(P\), coordinates can be selected that are as Cartesian as possible. The lecture initially uses the phrase Gaussian normal coordinates and then immediately qualifies it. The name is less important than the local conditions:
\begin{equation}
g_{mn}(P)=\delta_{mn},
\qquad
\partial_r g_{mn}(P)=0.
\label{eq:l4-normal-point}
\end{equation}
These are normal coordinates at one point. They need not remain Cartesian in a finite neighborhood, and many such frames exist because their axes may be rotated.

The metric is compatible with the derivative operator:
\begin{equation}
\nabla_r g_{mn}=0,
\qquad
\nabla_r g^{mn}=0.
\label{eq:l4-metric-compatibility}
\end{equation}
For a covector \(V_s\), the covariant derivative is
\begin{equation}
\nabla_r V_s
=
\partial_r V_s-\Gamma^p{}_{rs}V_p.
\label{eq:l4-covector-derivative}
\end{equation}
The connection coefficients are
\begin{equation}
\Gamma^p{}_{rs}
=
\frac{1}{2}g^{pn}
\left(
\partial_s g_{nr}
+\partial_r g_{ns}
-\partial_n g_{rs}
\right).
\label{eq:l4-christoffel}
\end{equation}
Several indices are corrected while this expression is written on the board. Equation \eqref{eq:l4-christoffel} is the completed relation. It is symmetric in \(r\) and \(s\), as expected for the torsion-free Levi-Civita connection.

\begin{figure}[htbp]
\centering
\includegraphics[width=0.92\textwidth]{lecture_04_figure_01.png}
\caption{The covector derivative and the completed Christoffel-symbol formula.}
\end{figure}

Because the Christoffel symbols are made from first derivatives of the metric, they vanish at \(P\) in the coordinates \eqref{eq:l4-normal-point}. They do not vanish in every coordinate system, even when the underlying space is flat. Thus \(\Gamma^p{}_{rs}\) is not a tensor. That non-tensorial character will later be exactly what allows an accelerated observer to describe a gravitational force that an inertial observer does not see.

\subsection{Raising the vector index}

The corresponding derivative of a vector is not introduced as a separate rule to memorize. It follows from the metric. Begin with
\begin{equation}
V^n=g^{np}V_p.
\end{equation}
Then
\begin{align}
\nabla_m V^n
&=\nabla_m\!\left(g^{np}V_p\right) \\
&=(\nabla_m g^{np})V_p+g^{np}\nabla_m V_p \\
&=g^{np}\left(\partial_mV_p-\Gamma^r{}_{mp}V_r\right),
\end{align}
where metric compatibility removed the first term. Reorganizing the indices gives
\begin{equation}
\nabla_m V^n
=
\partial_m V^n+\Gamma^n{}_{mr}V^r.
\label{eq:l4-vector-derivative}
\end{equation}

\begin{figure}[htbp]
\centering
\includegraphics[width=0.90\textwidth]{lecture_04_figure_02.png}
\caption{Raising the index with the inverse metric before differentiating.}
\end{figure}

\begin{figure}[htbp]
\centering
\includegraphics[width=0.90\textwidth]{lecture_04_figure_03.png}
\caption{The upper-index derivative emerging from metric compatibility and the product rule.}
\end{figure}

The index grammar is therefore simple:
\begin{equation}
\boxed{
\begin{aligned}
\text{lower index:}&\quad -\Gamma,\\
\text{upper index:}&\quad +\Gamma.
\end{aligned}}
\end{equation}
For a higher-rank tensor, one connection term appears for each index, with its sign determined by whether that index is lower or upper.

\subsection{A terminology question}

What exactly are the best coordinates mentioned above? They are not being asserted as a distinguished global coordinate system. At one selected point they make the metric equal to its flat form and make its first derivatives vanish. Rotating the axes gives another equally good frame. The geometric statements derived in such a frame do not depend on that orientation, because the final objects are tensors.

\section{Parallel Transport Along a Curve}

Suppose a vector \(V^n\) is carried along a curve. Over an infinitesimal displacement \(dx^m\), its covariant change is
\begin{align}
\delta V^n
&=(\nabla_mV^n)\,dx^m \\
&=\partial_mV^n\,dx^m
+\Gamma^n{}_{mr}V^r\,dx^m \\
&=dV^n+\Gamma^n{}_{mr}V^r\,dx^m.
\label{eq:l4-covariant-differential}
\end{align}

\begin{figure}[htbp]
\centering
\includegraphics[width=0.92\textwidth]{lecture_04_figure_04.png}
\caption{The ordinary component change and the connection correction along a curve.}
\end{figure}

The vector is parallel transported when this covariant change vanishes:
\begin{equation}
\boxed{
dV^n+\Gamma^n{}_{mr}V^r\,dx^m=0.
}
\label{eq:l4-parallel-transport}
\end{equation}
Operationally, we may imagine erecting a normal frame at each point of the curve. In that local frame, the transported vector is required not to turn. We then move an infinitesimal distance, erect the next normal frame, and repeat.

\begin{figure}[htbp]
\centering
\includegraphics[width=0.92\textwidth]{lecture_04_figure_05.png}
\caption{A clear board state of the parallel-transport equation.}
\end{figure}

The orientation of each temporary normal frame is not part of the physics. We could rotate it before making the comparison. What matters is that the covariant statement \eqref{eq:l4-parallel-transport} is independent of that arbitrary choice.

\subsection{Path dependence and curvature}

In a curved space, parallel transport depends on the path. Begin with the same vector at the same initial point and transport it to the same endpoint along two different routes. In general, the final vectors differ. Carrying a vector around a closed loop may return it rotated. That mismatch is holonomy, and it is a direct sign of curvature.

\begin{figure}[htbp]
\centering
\includegraphics[width=0.90\textwidth]{lecture_04_figure_06.png}
\caption{Two routes on a curved surface need not transport a vector to the same final orientation.}
\end{figure}

The cone provides a useful example. Away from its tip, a cone can be cut open and laid flat. Parallel transport around a loop that encloses the tip nevertheless returns the vector rotated by the angular deficit. If the tip is rounded, the curvature is finite and spread over a small region. If it is perfectly sharp, the idealization concentrates the curvature at the tip.

\begin{figure}[htbp]
\centering
\includegraphics[width=0.90\textwidth]{lecture_04_figure_07.png}
\caption{Circular motion on the cone and its cut-open planar representation.}
\end{figure}

\subsection{Questions about the transported vector}

\paragraph{Must the vector be defined everywhere?}
No. One may start with a single vector at one point and solve \eqref{eq:l4-parallel-transport} only along the chosen curve. A full vector field is optional. If a field has been defined on a region, it may be covariantly constant along one direction and not along another.

\paragraph{Is its length preserved?}
Yes. Let \(V^2=g_{mn}V^mV^n\). Along the curve,
\begin{align}
\frac{D}{d\lambda}\left(g_{mn}V^mV^n\right)
&=
\frac{Dg_{mn}}{d\lambda}V^mV^n
+2g_{mn}V^m\frac{DV^n}{d\lambda}\\
&=0.
\end{align}
Metric compatibility kills the first term and parallel transport kills the second. Lengths and mutual angles of parallel-transported vectors are therefore preserved.

\paragraph{Is the tangent to every curve parallel transported?}
No. This question leads directly to geodesics. The tangent to a circular route around a cone continually turns. The tangent to the straight line obtained after cutting the cone open does not. Only the second path is locally straight.

\section{Geodesics: Straightest Rather Than Merely Shortest}

A geodesic has two related descriptions. Globally, it is a stationary path of the length or proper-time functional. Locally, it is the straightest possible path. The word stationary is important: depending on the geometry and endpoints, the path need not be the unique global minimum. In Lorentzian spacetime, timelike geodesics locally maximize proper time rather than minimize a positive distance.

The local description is the one needed here. A small car driven over curved terrain with its steering wheel held dead ahead traces a geodesic, provided the car is small compared with the scale on which the terrain bends. It need not look straight in an embedding space; it is straight according to the geometry available on the surface.

\begin{figure}[htbp]
\centering
\includegraphics[width=0.90\textwidth]{lecture_04_figure_08.png}
\caption{A path drawn intrinsically on the curved surface, preparing the straightest-path definition.}
\end{figure}

\subsection{The tangent vector}

Let the curve be \(x^m(s)\), with arc length
\begin{equation}
ds^2=g_{mn}\,dx^m dx^n.
\end{equation}
The unit tangent is
\begin{equation}
t^m=\frac{dx^m}{ds}.
\label{eq:l4-unit-tangent}
\end{equation}
Geometrically, take two nearby points, form the displacement vector between them, divide by their separation, and let the points merge.

\begin{figure}[htbp]
\centering
\includegraphics[width=0.90\textwidth]{lecture_04_figure_09.png}
\caption{The tangent direction as the limiting displacement along a curve.}
\end{figure}

\begin{figure}[htbp]
\centering
\includegraphics[width=0.90\textwidth]{lecture_04_figure_10.png}
\caption{The line element and the unit tangent \(t^m=dx^m/ds\).}
\end{figure}

A geodesic is a curve whose tangent is parallel transported along itself:
\begin{equation}
dt^n+\Gamma^n{}_{mr}t^r\,dx^m=0.
\end{equation}
Divide by \(ds\) and use \(dx^m/ds=t^m\):
\begin{equation}
\frac{dt^n}{ds}
=
-\Gamma^n{}_{mr}t^m t^r.
\end{equation}
Finally insert \(t^n=dx^n/ds\):
\begin{equation}
\boxed{
\frac{d^2x^n}{ds^2}
+\Gamma^n{}_{mr}
\frac{dx^m}{ds}
\frac{dx^r}{ds}
=0.
}
\label{eq:l4-geodesic}
\end{equation}

\begin{figure}[htbp]
\centering
\includegraphics[width=0.94\textwidth]{lecture_04_figure_11.png}
\caption{The tangent-vector transport law and the second-order geodesic equation.}
\end{figure}

Equation \eqref{eq:l4-geodesic} already resembles Newtonian dynamics. If \(s\) is replaced by an appropriate time parameter, the first term is an acceleration. The connection term then acts like a force per unit mass, determined entirely by the metric. In general relativity this analogy becomes literal: free particles follow spacetime geodesics, and what Newton called gravitational acceleration appears through the connection.

\section{From Riemannian Distance to Spacetime Proper Time}

The discussion now moves from curved space to spacetime. In inertial Cartesian coordinates, the proper time between nearby events is
\begin{equation}
d\tau^2
=dt^2-dx^2-dy^2-dz^2.
\label{eq:l4-proper-time}
\end{equation}
The lecture uses
\begin{equation}
d\tau^2=-g_{\mu\nu}\,dx^\mu dx^\nu
\label{eq:l4-lecture-sign}
\end{equation}
with
\begin{equation}
g_{\mu\nu}
=
\eta_{\mu\nu}
=
\begin{pmatrix}
-1&0&0&0\\
0&1&0&0\\
0&0&1&0\\
0&0&0&1
\end{pmatrix}.
\label{eq:l4-eta}
\end{equation}
Thus the external minus sign in \eqref{eq:l4-lecture-sign} converts the chosen mostly-plus metric into \eqref{eq:l4-proper-time}. The coordinates are
\begin{equation}
\begin{aligned}
x^\mu&=(t,x,y,z),\\
x^0&=t,\quad x^1=x,\\
x^2&=y,\quad x^3=z.
\end{aligned}
\end{equation}

\begin{figure}[htbp]
\centering
\includegraphics[width=0.94\textwidth]{lecture_04_figure_12.png}
\caption{Proper time, spacetime coordinates, and the mostly-plus Minkowski metric.}
\end{figure}

The signature is the invariant information: one eigenvalue has the sign associated with time and three have the sign associated with space. Multiplying the entire metric by \(-1\) and adjusting every convention consistently gives the equally common mostly-minus convention. Nothing physical depends on that overall choice.

The lecture asks whether spaces with other indefinite signatures have been considered. Mathematically, pseudo-Riemannian geometries with many signatures are well established. The physical point is narrower: observed relativistic spacetime is Lorentzian, with one temporal direction. Additional timelike directions would profoundly alter causal structure and do not describe the physics under discussion.

\begin{figure}[htbp]
\centering
\includegraphics[width=0.94\textwidth]{lecture_04_figure_13.png}
\caption{The sign convention is checked before the flat metric is denoted by \(\eta_{\mu\nu}\).}
\end{figure}

Flat spacetime does not mean that the metric components must look like \(\eta_{\mu\nu}\) in every coordinate system. It means that some coordinates exist in which they do. Curvilinear or accelerated coordinates can make the components vary with position even when the Riemann tensor is zero.

\section{Why Equal Coordinate Acceleration Is Not Rigid}

We want coordinates attached to an accelerating elevator. The first idea is to line up a set of observers at fixed separations and give them all the same coordinate acceleration. Relativity prevents that naive construction from remaining rigid.

\begin{figure}[htbp]
\centering
\includegraphics[width=0.92\textwidth]{lecture_04_figure_14.png}
\caption{A line of neighboring observers used to test whether a common acceleration preserves rigidity.}
\end{figure}

If separated observers follow identical acceleration programs according to one inertial frame, the distance between them in their changing instantaneous rest frames does not remain fixed. A string joining them stretches and eventually breaks. This is the same kinematic tension exhibited by the Bell-spaceship thought experiment. Moreover, constant Newtonian coordinate acceleration cannot continue indefinitely: it would eventually predict a speed greater than light.

A relativistically rigid accelerated frame must therefore use a different proper acceleration at each position. The observers nearer the would-be horizon accelerate more strongly; those farther away accelerate more gently. Their proper separations can then remain constant in the momentarily comoving frame. This is Born-rigid acceleration in one spatial dimension.

\section{Euclidean Polar Coordinates as the Warm-Up}

The correct construction is easiest to see by beginning with an ordinary circle. In a Euclidean plane,
\begin{equation}
x=r\cos\theta,
\qquad
y=r\sin\theta,
\end{equation}
so
\begin{equation}
x^2+y^2=r^2,
\qquad
\cos^2\theta+\sin^2\theta=1.
\end{equation}
The complex-exponential forms are
\begin{equation}
\cos\theta=\frac{e^{i\theta}+e^{-i\theta}}{2},
\qquad
\sin\theta=\frac{e^{i\theta}-e^{-i\theta}}{2i}.
\label{eq:l4-euler-trig}
\end{equation}
The factor of \(i\) in the denominator of the sine formula is caught and restored during the board discussion.

\begin{figure}[htbp]
\centering
\includegraphics[width=0.94\textwidth]{lecture_04_figure_15.png}
\caption{The circle, Euclidean polar coordinates, and their exponential trigonometric forms.}
\end{figure}

A fixed value of \(r\) is a circle. Motion with constant angular speed has constant centripetal acceleration in magnitude, always directed inward. The hyperbola will play the corresponding role in spacetime.

\section{Hyperbolic Coordinates in Minkowski Spacetime}

Replace trigonometric functions by their hyperbolic relatives:
\begin{equation}
\cosh\omega=\frac{e^\omega+e^{-\omega}}{2},
\qquad
\sinh\omega=\frac{e^\omega-e^{-\omega}}{2},
\end{equation}
with
\begin{equation}
\cosh^2\omega-\sinh^2\omega=1.
\label{eq:l4-hyperbolic-identity}
\end{equation}
There is no factor of \(i\). The coordinate \(\omega\) is a hyperbolic angle, or rapidity. Unlike an ordinary angle, it is not periodic; it ranges from \(-\infty\) to \(+\infty\).

\begin{figure}[htbp]
\centering
\includegraphics[width=0.94\textwidth]{lecture_04_figure_16.png}
\caption{Circular functions are replaced by hyperbolic functions in the spacetime construction.}
\end{figure}

Define new coordinates in the right Rindler wedge:
\begin{equation}
x=r\cosh\omega,
\qquad
t=r\sinh\omega,
\qquad
r>0.
\label{eq:l4-rindler-transform}
\end{equation}
Then
\begin{equation}
x^2-t^2=r^2.
\label{eq:l4-rindler-hyperbola}
\end{equation}
At fixed \(r\), varying \(\omega\) traces a hyperbola. Increasing \(r\) moves across the family in a spacelike direction, while increasing \(\omega\) moves forward along one hyperbola in a timelike direction.

\begin{figure}[htbp]
\centering
\includegraphics[width=0.94\textwidth]{lecture_04_figure_17.png}
\caption{The coordinate transformation \(x=r\cosh\omega\), \(t=r\sinh\omega\), and its invariant hyperbola.}
\end{figure}

\begin{figure}[htbp]
\centering
\includegraphics[width=0.94\textwidth]{lecture_04_figure_18.png}
\caption{A congruence of fixed-\(r\) hyperbolae representing a rigidly accelerated family.}
\end{figure}

\subsection{The metric in accelerated coordinates}

Differentiate \eqref{eq:l4-rindler-transform}:
\begin{align}
dx&=\cosh\omega\,dr+r\sinh\omega\,d\omega,\\
dt&=\sinh\omega\,dr+r\cosh\omega\,d\omega.
\end{align}
Using \eqref{eq:l4-hyperbolic-identity}, the cross terms cancel:
\begin{align}
d\tau^2
&=dt^2-dx^2\\
&=r^2\,d\omega^2-dr^2.
\label{eq:l4-rindler-metric}
\end{align}
The transverse directions, when restored, simply add \(-dy^2-dz^2\) to the proper-time convention used here.

This is still flat spacetime. The coordinate-dependent coefficient \(r^2\) is no more evidence of curvature than the \(r^2\) multiplying \(d\theta^2\) in Euclidean polar coordinates.

\subsection{Proper acceleration along the hyperbolae}

For a fixed worldline \(r=R\), equation \eqref{eq:l4-rindler-metric} gives
\begin{equation}
d\tau=R\,d\omega.
\end{equation}
Its four-velocity is
\begin{equation}
U^\mu
=
\frac{d(t,x)}{d\tau}
=
(\cosh\omega,\sinh\omega),
\end{equation}
and its four-acceleration is
\begin{equation}
A^\mu
=
\frac{dU^\mu}{d\tau}
=
\frac{1}{R}(\sinh\omega,\cosh\omega).
\end{equation}
The invariant magnitude is constant:
\begin{equation}
a=\sqrt{A^\mu A_\mu}=\frac{1}{R}
\qquad (c=1).
\label{eq:l4-proper-acceleration}
\end{equation}

\begin{figure}[htbp]
\centering
\includegraphics[width=0.90\textwidth]{lecture_04_figure_19.png}
\caption{The inverse-radius law \(a=1/R\) for a fixed accelerated worldline.}
\end{figure}

With units restored,
\begin{equation}
\boxed{a=\frac{c^2}{R}.}
\label{eq:l4-proper-acceleration-units}
\end{equation}

\begin{figure}[htbp]
\centering
\includegraphics[width=0.90\textwidth]{lecture_04_figure_20.png}
\caption{Restoring units converts \(1/R\) into \(c^2/R\).}
\end{figure}

The smaller the intercept \(R\), the more sharply the worldline bends and the larger the proper acceleration felt by its observer. Farther hyperbolae bend less and require less acceleration. Each observer has constant acceleration in time, but different observers have different accelerations. That variation is precisely what preserves their mutual proper distances.

\subsection{Questions about the accelerated coordinates}

\paragraph{Is \(R\) the \(x\)-intercept?}
Yes. At \(\omega=0\), equation \eqref{eq:l4-rindler-transform} gives \(t=0\) and \(x=R\). The intercept fixes the proper acceleration through \eqref{eq:l4-proper-acceleration-units}.

\paragraph{Does the construction depend on the chosen origin?}
The coordinates do. Translating the origin produces another family of hyperbolae with shifted asymptotes. No invariant physics depends on that choice. We are deliberately selecting one accelerated coordinate system and asking how inertial motion appears when described in it.

\paragraph{What is proper acceleration?}
It is the acceleration measured by an accelerometer carried with the observer. Covariantly,
\begin{equation}
A^\mu=\frac{dU^\mu}{d\tau},
\qquad
U^\mu=\frac{dx^\mu}{d\tau}.
\end{equation}
Its magnitude is the push the observer actually feels, not the coordinate acceleration \(d^2x/dt^2\) assigned by a distant frame.

\section{Flat Metrics in Polar and Hyperbolic Coordinates}

Before localizing near the Earth, it is useful to place the Euclidean and Lorentzian metrics side by side. In the plane,
\begin{equation}
ds^2=dr^2+r^2d\theta^2,
\qquad
g_{ab}
=
\begin{pmatrix}
1&0\\
0&r^2
\end{pmatrix}.
\label{eq:l4-polar-metric}
\end{equation}
An audience suggestion briefly introduces a spherical factor, but this is only a plane, not a sphere. The \(r^2\) term records the fact that a fixed angular interval spans a longer arc at larger radius.

\begin{figure}[htbp]
\centering
\includegraphics[width=0.94\textwidth]{lecture_04_figure_21.png}
\caption{The Euclidean polar metric: position-dependent components in a flat plane.}
\end{figure}

The Minkowski analog is
\begin{equation}
d\tau^2=r^2d\omega^2-dr^2.
\end{equation}
Both metrics have position-dependent components, yet both geometries are flat. Coordinate curvature and geometric curvature are not the same thing.

\section{Localizing Near Terrestrial Acceleration}

Choose one reference hyperbola whose proper acceleration is the terrestrial value \(g\):
\begin{equation}
g=\frac{c^2}{R},
\qquad
R=\frac{c^2}{g}.
\label{eq:l4-earth-radius}
\end{equation}

\begin{figure}[htbp]
\centering
\includegraphics[width=0.94\textwidth]{lecture_04_figure_22.png}
\caption{Selecting the reference hyperbola by \(c^2/R=g\).}
\end{figure}

Numerically,
\begin{equation}
R
\simeq
\frac{(3.0\times10^8\ {\rm m\,s^{-1}})^2}
{9.8\ {\rm m\,s^{-2}}}
\simeq
9.2\times10^{15}\ {\rm m}.
\end{equation}
The lecture rounds this to \(10^{16}\) meters. A field of one terrestrial \(g\) is modest on relativistic scales, so the corresponding Rindler radius is enormous.

Now examine only a small neighborhood of that reference observer. Introduce a local height \(y\) by
\begin{equation}
r=R+y,
\qquad
dr=dy,
\qquad
|y|\ll R.
\label{eq:l4-local-height}
\end{equation}
The equality \(dr=dy\) is explicitly caught during the discussion, so the spatial term must be written as \(-dy^2\).

Substitution into \eqref{eq:l4-rindler-metric} gives
\begin{align}
d\tau^2
&=(R+y)^2d\omega^2-dy^2\\
&=(R^2+2Ry+y^2)d\omega^2-dy^2\\
&=
\left(
1+\frac{2y}{R}+\frac{y^2}{R^2}
\right)
R^2d\omega^2-dy^2.
\label{eq:l4-local-expand}
\end{align}

\begin{figure}[htbp]
\centering
\includegraphics[width=0.94\textwidth]{lecture_04_figure_23.png}
\caption{Expanding the accelerated metric about \(r=R+y\).}
\end{figure}

Define a local time coordinate
\begin{equation}
t=R\omega.
\label{eq:l4-local-time}
\end{equation}
Then
\begin{equation}
d\tau^2
=
\left(
1+\frac{2y}{R}+\frac{y^2}{R^2}
\right)dt^2-dy^2.
\label{eq:l4-local-exact}
\end{equation}

\begin{figure}[htbp]
\centering
\includegraphics[width=0.94\textwidth]{lecture_04_figure_24.png}
\caption{Introducing \(t=R\omega\) after the expansion about the reference hyperbola.}
\end{figure}

Even a vertical excursion of tens of kilometers is tiny compared with \(R\). The lecture uses a fall of roughly twenty-four miles to make the comparison vivid: such a distance is still negligible beside \(10^{16}\) meters. Therefore \(y^2/R^2\) is much smaller than \(y/R\), and the quadratic term may be dropped:
\begin{equation}
d\tau^2
\simeq
\left(1+\frac{2y}{R}\right)dt^2-dy^2.
\label{eq:l4-local-linear}
\end{equation}

\begin{figure}[htbp]
\centering
\includegraphics[width=0.94\textwidth]{lecture_04_figure_25.png}
\caption{The expanded metric and the local-time substitution before the weak-field truncation.}
\end{figure}

In units with \(c=1\), equation \eqref{eq:l4-earth-radius} gives \(1/R=g\), and hence
\begin{equation}
\boxed{
d\tau^2\simeq(1+2gy)\,dt^2-dy^2.
}
\label{eq:l4-uniform-metric}
\end{equation}

\begin{figure}[htbp]
\centering
\includegraphics[width=0.94\textwidth]{lecture_04_figure_26.png}
\caption{The completed local metric for the uniformly accelerated elevator.}
\end{figure}

The quantity
\begin{equation}
\Phi(y)=gy
\end{equation}
is the Newtonian gravitational potential per unit mass in a uniform field. Thus
\begin{equation}
d\tau^2\simeq(1+2\Phi)\,dt^2-dy^2
\qquad (c=1).
\label{eq:l4-weak-potential-propertime}
\end{equation}
With the lecture's convention \(d\tau^2=-g_{\mu\nu}dx^\mu dx^\nu\), this means
\begin{equation}
g_{tt}=-(1+2\Phi).
\label{eq:l4-gtt-potential}
\end{equation}
In conventional units the dimensionless correction is \(2\Phi/c^2\).

Nothing has curved spacetime. Equation \eqref{eq:l4-uniform-metric} is the local expansion of a flat Rindler metric. Its apparent gravity is the inertial force seen inside an accelerated elevator.

\section{The Geodesic Equation Becomes Newton's Law}

The identification of \(gy\) with a potential has so far been suggestive. The decisive test is dynamical: does a free particle in the metric \eqref{eq:l4-uniform-metric} accelerate downward with \(d^2y/dt^2=-g\)?

For a timelike worldline parametrized by proper time,
\begin{equation}
\frac{d^2x^\nu}{d\tau^2}
+\Gamma^\nu{}_{\mu\rho}
\frac{dx^\mu}{d\tau}
\frac{dx^\rho}{d\tau}
=0.
\label{eq:l4-spacetime-geodesic}
\end{equation}
The vertical component is
\begin{equation}
\frac{d^2y}{d\tau^2}
=
-\Gamma^y{}_{\mu\rho}
\frac{dx^\mu}{d\tau}
\frac{dx^\rho}{d\tau}.
\label{eq:l4-y-geodesic}
\end{equation}

\begin{figure}[htbp]
\centering
\includegraphics[width=0.94\textwidth]{lecture_04_figure_27.png}
\caption{The transport equation and its second-order geodesic form, ready to be applied to vertical motion.}
\end{figure}

\subsection{The slow-motion approximation}

For a slowly moving particle in a weak, static field,
\begin{equation}
\frac{dt}{d\tau}\simeq1,
\qquad
\left|\frac{dy}{d\tau}\right|\ll1.
\end{equation}
The time component of the four-velocity is therefore of order one, while the spatial components are small compared with it. Among the products in \eqref{eq:l4-y-geodesic}, the dominant one has \(\mu=\rho=t\):
\begin{equation}
\frac{d^2y}{d\tau^2}
\simeq
-\Gamma^y{}_{tt}
\left(\frac{dt}{d\tau}\right)^2
\simeq
-\Gamma^y{}_{tt}.
\label{eq:l4-slow-geodesic}
\end{equation}

\begin{figure}[htbp]
\centering
\includegraphics[width=0.94\textwidth]{lecture_04_figure_28.png}
\caption{Specializing the spacetime geodesic equation to the slowly moving \(y\)-component.}
\end{figure}

Now calculate the one required connection coefficient. The metric is diagonal and independent of \(t\), so
\begin{align}
\Gamma^y{}_{tt}
&=
\frac{1}{2}g^{yy}
\left(
\partial_tg_{yt}
+\partial_tg_{ty}
-\partial_yg_{tt}
\right)\\
&=
-\frac{1}{2}\partial_yg_{tt},
\label{eq:l4-gamma-ytt}
\end{align}
because \(g^{yy}=1\) in the lecture's mostly-plus tensor convention.

\begin{figure}[htbp]
\centering
\includegraphics[width=0.94\textwidth]{lecture_04_figure_29.png}
\caption{The Christoffel formula reduced to the \(y,t,t\) component.}
\end{figure}

Combining \eqref{eq:l4-slow-geodesic} and \eqref{eq:l4-gamma-ytt},
\begin{equation}
\frac{d^2y}{d\tau^2}
\simeq
\frac{1}{2}\partial_yg_{tt}.
\label{eq:l4-accel-gtt}
\end{equation}
From \eqref{eq:l4-gtt-potential},
\begin{equation}
\partial_yg_{tt}
=
-2\partial_y\Phi
=
-2g.
\end{equation}
Therefore
\begin{equation}
\boxed{
\frac{d^2y}{d\tau^2}=-g.
}
\label{eq:l4-newton-result}
\end{equation}
Since \(t\simeq\tau\) in the same slow-motion approximation, this is also
\begin{equation}
\frac{d^2y}{dt^2}=-g.
\end{equation}

\begin{figure}[htbp]
\centering
\includegraphics[width=0.94\textwidth]{lecture_04_figure_30.png}
\caption{The completed sign chain: the \(y\)-derivative of \(g_{tt}\) produces the acceleration \(-g\).}
\end{figure}

More generally, in the Newtonian regime,
\begin{equation}
\boxed{
\begin{aligned}
g_{tt}&=-(1+2\Phi),\\
\frac{d^2x^i}{dt^2}&=-\partial_i\Phi
\end{aligned}
}
\qquad(c=1).
\label{eq:l4-newton-general}
\end{equation}
With units restored, the metric correction is \(2\Phi/c^2\), while the acceleration remains \(-\boldsymbol{\nabla}\Phi\).

This calculation explains why the metric, rather than a separate gravitational force field, is the central object. A free particle simply follows a spacetime geodesic. In coordinates adapted to an accelerated laboratory, that geodesic has the coordinate curvature ordinarily attributed to gravity.

\subsection{A question about the chosen Rindler origin}

Could the reference hyperbola have been shifted so that its asymptotes and intercept occur elsewhere? Yes. There is a uniformly accelerated coordinate system associated with every suitable translated origin. In the elevator picture, these choices amount to selecting elevators whose floors occupy different positions at one instant.

The local result does not rely on a privileged origin. Once a reference observer with acceleration \(g\) has been selected, \(y\) measures height relative to that observer and the nearby metric takes the form \eqref{eq:l4-uniform-metric}. The arbitrary coordinate choice determines where we describe the floor; it does not create invariant curvature.

\section{Why the Christoffel Symbol Must Not Be a Tensor}

The post-lecture discussion returns to a point that first appeared during the review. In the original inertial coordinates,
\begin{equation}
g_{\mu\nu}=\eta_{\mu\nu},
\qquad
\partial_\rho\eta_{\mu\nu}=0,
\qquad
\Gamma^\mu{}_{\nu\rho}=0.
\end{equation}
A free particle moves on a straight line, and there is no gravitational force.

In the accelerated coordinates, the same flat spacetime has position-dependent metric components. Its Christoffel symbols do not vanish, and the same inertial trajectory appears curved. The particle seems to fall toward the floor of the elevator.

This would be impossible if \(\Gamma^\mu{}_{\nu\rho}\) were a tensor. A tensor that vanishes in one frame vanishes in every frame. The connection instead acquires an inhomogeneous term under a coordinate transformation, so it may be zero in an inertial frame and nonzero in an accelerated one. That is exactly what the equivalence principle requires:
\begin{equation}
\text{accelerated frame}
\quad\longleftrightarrow\quad
\text{local gravity}.
\end{equation}

Curvature behaves differently. The Riemann tensor is a tensor. If it vanishes in the inertial frame, it vanishes in the accelerated frame as well:
\begin{equation}
R^\mu{}_{\nu\rho\sigma}=0.
\end{equation}
Thus the Christoffel symbols describe coordinate-dependent gravitational force, while curvature detects the invariant tidal content that no change of coordinates can remove.

\section{From a Uniform Potential to a Schwarzschild Preview}

The accelerated elevator has no gravitating matter and no curvature. Nevertheless, the weak-field relation \eqref{eq:l4-newton-general} suggests what to look for in the field of an actual spherical mass. The Newtonian potential outside a mass \(M\) is
\begin{equation}
\Phi(r)=-\frac{GM}{r}.
\label{eq:l4-newton-potential}
\end{equation}
Substituting it into the weak-field time component gives
\begin{equation}
g_{tt}
\simeq
-\left(1-\frac{2GM}{r}\right)
\qquad(c=1).
\label{eq:l4-schwarzschild-preview}
\end{equation}
The sign of the potential is corrected during the board discussion. Attraction requires a negative potential that approaches zero from below as \(r\) tends to infinity.

\begin{figure}[htbp]
\centering
\includegraphics[width=0.94\textwidth]{lecture_04_figure_31.png}
\caption{The Newtonian potential inserted into the time-time metric coefficient, revealing the factor \(1-2GM/r\).}
\end{figure}

The coefficient vanishes at
\begin{equation}
r=2GM
\qquad(c=1),
\end{equation}
or, with units restored,
\begin{equation}
r_s=\frac{2GM}{c^2}.
\end{equation}
This is the Schwarzschild radius. The vanishing of this particular metric coefficient foreshadows a horizon.

The argument is deliberately presented as a preview, not a derivation. The weak-field approximation alone cannot establish the exact Schwarzschild solution near \(r=2GM/c^2\), precisely where the correction is no longer small. The full solution also contains a nontrivial radial component and must satisfy Einstein's field equations. Those equations have not yet been introduced. What the present reasoning reveals is the correct structural clue and the scale at which something qualitatively new must occur.

\section{Classroom Questions and Corrections}

\subsection{Does a hyperbolic angle behave like an ordinary angle?}

No. Both parameterize motion along a family of constant-radius curves, but \(\theta\) is periodic while the rapidity \(\omega\) ranges over the real line. The word angle refers to the algebraic analogy between rotations and Lorentz boosts, not to a quantity restricted between \(0\) and \(2\pi\).

\subsection{Why does a modest acceleration require an enormous R?}

The relation \(a=c^2/R\) contains the large relativistic scale \(c^2\). To obtain only \(10\ {\rm m\,s^{-2}}\), one needs \(R\) of order \(10^{16}\) meters. This is not a new physical distance attached to the Earth. It is the curvature radius of the hyperbolic worldline in a Minkowski diagram, expressed in length units.

The large radius also guarantees a useful Newtonian interval. Over laboratory times and distances, the speed remains far below \(c\), the acceleration varies negligibly across the apparatus, and \(t\) differs little from proper time.

\subsection{Is the apparent gravitational field real?}

It is real as an accelerometer reading for an observer held at fixed Rindler position, but it is not invariant curvature. Release the observer, and the accelerometer reads zero along the resulting inertial geodesic. A sufficiently small freely falling laboratory eliminates the uniform field locally. What cannot be eliminated over a finite region is a tidal field, represented by nonzero curvature.

\subsection{What was corrected during the force calculation?}

Several signs and a factor of one half are checked live. The reliable chain in the lecture's convention is
\begin{equation}
\begin{aligned}
g_{tt}&=-(1+2gy),\\
\Gamma^y{}_{tt}&=-\frac12\partial_yg_{tt}=g,\\
\frac{d^2y}{d\tau^2}&=-\Gamma^y{}_{tt}=-g.
\end{aligned}
\end{equation}
Writing the convention explicitly prevents a board-level sign hesitation from becoming a physical ambiguity.

\section{What Has Been Established}

The lecture starts with differential geometry and ends with Newton's falling body. The essential steps are:
\begin{enumerate}
\item A covariant derivative compares tensors in changing coordinate bases.
\item Parallel transport means zero covariant change along a chosen path.
\item A geodesic parallel transports its own tangent.
\item Free particles follow geodesics of spacetime.
\item Fixed-\(r\) hyperbolae form a relativistically rigid accelerated frame.
\item Their proper acceleration is \(a=c^2/r\), so it must vary with position.
\item Near a reference worldline with acceleration \(g\), the flat metric is
\begin{equation}
d\tau^2\simeq(1+2gy)\,dt^2-dy^2
\qquad(c=1).
\end{equation}
\item Its geodesic equation reduces to
\begin{equation}
\frac{d^2y}{dt^2}=-g.
\end{equation}
\end{enumerate}

The equivalence principle is therefore not merely an analogy between an elevator and gravity. It is encoded in the way the metric and its connection transform. A uniform gravitational force can appear when flat spacetime is written in accelerated coordinates, while the curvature remains zero.

The final extrapolation points beyond the uniform field. Replacing \(\Phi=gy\) by \(\Phi=-GM/r\) predicts the characteristic factor \(1-2GM/r\) in the time component of a spherical metric and exposes the Schwarzschild scale. To decide the full geometry, however, we need the field equations. That is the task to which the next lecture turns.
